\section{TapeAgents: trazas de ejecución auditables}
Chapados presenta TapeAgents entre las 00:18:56 y las 00:20:52; se trata de un marco de trabajo recién liberado como código abierto cuyo objetivo es «tender un puente entre las herramientas de ingeniería y las herramientas de optimización». «Tape» es la abstracción unificada: en este registro se anotan el razonamiento, las acciones y las observaciones de todos los Agentes, por lo que resulta auditable y además puede ser leído por Agentes posteriores.

\begin{importantbox}{El «Best of Both Worlds» de TapeAgents}
«Queremos herramientas de desarrollo y depuración, pero también capacidad de optimización de prompts y de fine-tuning»; TapeAgents ofrece componentes de grano fino, flujos paralelos y runtime orchestration, y a la vez usa el tape como activo de datos para convertir el registro de ejecución en material de optimización.
\end{importantbox}

\subsection{Estructura y reutilización del tape}
El tape de TapeAgents no es un simple registro: contiene metadata estructurada y abundante que puede ser consumida por otros Agentes o por optimizadores (00:21:10--00:22:25). Varios Agentes pueden compartir el mismo tape; el Agente B lee el historial completo y elige el node que desea ejecutar, con lo que se construye un flujo de control modular.

\begin{knowledgebox}{El tape como «datos como servicio»}
El tape registra cada prompt, cada retroalimentación del entorno y cada instrucción de acción, lo que permite impulsar la depuración, la visualización, la optimización de prompts y la destilación maestro-alumno, formando así un ciclo cerrado que va del execution trace a los datos de fine-tune.
\end{knowledgebox}

\subsection{Resumen de la sección}
Gracias a un tape unificado con metadata abundante, TapeAgents logra la sinergia entre la colaboración entre varios Agentes, la depuración a nivel de ingeniería y la optimización automatizada, resolviendo la desconexión entre el «prompt manual» y la «optimización experimental».

